\documentclass[12pt]{article}
\usepackage{natbib}
\usepackage{floatrow,threeparttable}
\floatsetup[figure]{capposition=top}
\floatsetup[table]{capposition=top}
\usepackage{ragged2e}
\usepackage[left=1in, right=1in,top=1.25in,bottom=1.25in]{geometry}
\usepackage{multirow}
\usepackage{bbm}
\usepackage{lscape}
%\usepackage[title,titletoc,toc]{appendix}
\usepackage{xparse}
%\usepackage[bf,small]{titlesec}
\usepackage{xr}
\usepackage{amsmath,amssymb,amsthm,mathrsfs,amsfonts,dsfont}
\usepackage{commath}
\usepackage{graphicx} 
\usepackage{booktabs}
\usepackage[bf]{titlesec}
\usepackage{setspace}
\usepackage{pdflscape}

\usepackage{bibunits}

\usepackage[nottoc,numbib]{tocbibind}

\defaultbibliography{bibliography.bib}
\defaultbibliographystyle{ecta}


\usepackage{verbatim}
%\usepackage[large]{subfigure}
\usepackage{subfigure}
\bibliographystyle{chicago}
\usepackage{soul}
\usepackage[addedmarkup=colored,deletedmarkup=sout,commandnameprefix=always]{changes}
\usepackage{tikz}
\usetikzlibrary{datavisualization}
\usetikzlibrary{datavisualization.formats.functions}
\usetikzlibrary{decorations.pathreplacing}
\tikzset{every picture/.style={line width=0.75pt}} 
\usepackage{colortbl}
\usepackage[bf]{caption}
\usepackage{enumerate}
\usepackage[bookmarks, pdftitle={}, pdfauthor={}, colorlinks=true, linkcolor=red, citecolor=blue, urlcolor=blue]{hyperref}

%\usepackage{booktabs}
%\usepackage{times} 
\usepackage{txfonts}
\usetikzlibrary{decorations.pathreplacing}
\usepackage[multiple,stable]{footmisc}
\setlength{\footnotesep}{.1in}
\linespread{1.25}
\setlength{\parskip}{6pt}
\usepackage{datetime}
\newdateformat{monthyeardate}{\monthname[\THEMONTH], \THEYEAR}

\begin{document}



%%%%%%%%%%%%%%%%%%%%%%%%%%%%%%%%%%%%%%%%%
%% 					TITLE								       %%
%%%%%%%%%%%%%%%%%%%%%%%%%%%%%%%%%%%%%%%%%


\title{\Large {Do CCT Programs (Really) Reduce Mortality:  10 Year Evidence from Mexico}\thanks{This version: March 2026. The opinions expressed in this paper are those of the authors and do not necessarily reflect the views of the Inter-American Development Bank, its Board of Directors, or the countries they represent, nor of INEGI. All remaining errors are our own. All data used in this study are publicly available from INEGI.}}

\author{Emma Aguila\footnote{Emma Aguila is an Associate Professor in the Sol Price School of Public Policy at the University of Southern California} \vspace{-4mm} \and William H. Dow\footnote{William H. Dow is a Professor of Health Policy and Management in the School of Public Health, as well as Professor in the Department of Demography at the University of California, Berkeley.}  \and Felipe Menares\footnote{Felipe Menares is an Economic/Impact Evaluation Consultant at the Inter-American Development Bank. Email: \texttt{felipeme@iadb.org}} \and Susan W. Parker \footnote{Susan W. Parker is Professor of Public Policy in the School of Public Policy at the University of Maryland }  \and Jorge Peniche\footnote{Jorge Peniche is a Researcher in the Sol Price School of Public Policy at the University of Southern California} \and Soomin Ryu\footnote{Soomin Ryu is an Assistant Professor in the Department of Public Health at the University of Missouri}}


\date{}

\maketitle

\vspace{-12mm}

\thispagestyle{empty}

%%%%%%%%%%%%%%%%%%%%%%%%%%%%%%%%%%%%%%%%%
%% 					ABSTRACT							       %%
%%%%%%%%%%%%%%%%%%%%%%%%%%%%%%%%%%%%%%%%%

\begin{abstract}
%\singlespacing

\noindent 
Early work on the impacts of the pioneering program Progresa suggested large, short-term effects of the program to reduce mortality of older adults in Mexico. We re-examine the effects of Progresa on elderly mortality rates after 10 years and find little evidence of lasting impacts of the Progresa program in reducing elderly mortality. We employ difference-in-differences models that exploit geographic variation in program expansion to estimate effects, focusing on high-poverty municipalities.  We argue these results may not be surprising. Aging individuals were not generally the recipients in conditional cash transfer programs and the benefits provided were largely conditioned to school enrollment of children in the household. Health conditionalities required aging individuals to attend health clinics only once per year, a relatively minimal requirement. We re-analyze prior evidence suggesting large negative short-term impacts on mortality, and find that these results are somewhat sensitive, becoming smaller and/or insignificant with relatively minor changes in the specification.  Combined with our own estimates, we conclude there is little evidence that Progresa reduced aging mortality during its first 10 years of operation. 




\noindent  \\
\vspace{0in}\\
\noindent\textbf{Keywords: Mortality, Poverty, Demography} \\
\vspace{0in}\\
\noindent\textbf{JEL Codes: I13, I14, I15, I18} 


\end{abstract}
%%%%%%%%%%%%%%%%%%%%%%%%%%%%%%%%%%%%%%%%%
%% 					INTRO								       %%
%%%%%%%%%%%%%%%%%%%%%%%%%%%%%%%%%%%%%%%%%
\clearpage

\setcounter{page}{1}

\section{Introduction}
Research on the mortality effects of income-support social insurance programs for older adults has generated conflicting results. Some studies suggest important health benefits, others find no effects, and still others find unintended adverse effects potentially linked to pathways such as increased obesity. Furthermore, most evidence has focused predominantly on short-run effects rather than net long-run mortality effects.  

Mexico offers an important opportunity to study the longer-term effects of social insurance programs on aging mortality. The Progresa conditional cash transfer (CCT) program, which began in 1997, was renamed  Oportunidades in 2001 under the Fox administration, and then renamed as Prospera in 2007 under the Calderon administration. Progresa has drawn considerable attention from researchers and policymakers \citep{LevySchady2013JEP}, and has served as a model for similar programs in more than 60 countries \citep{Baird2014JDEffectiveness, Fiszbein2009CCTBook, ParkerTodd2017JEL}. There are numerous studies on the short-term health effects of Progresa \citep{Gertler2004AER, LagardeHainesPalmer2009Cochrane, ParkerTodd2017JEL, RiveraHernandez2016HSR}, but little evidence on its long-term effects \citep{ParkerTodd2017JEL}.  
    
Furthermore, although conditional cash transfer (CCT) programs are widespread, there is remarkably little evidence regarding their effects on well-being and health of older adults.  Much research has emphasized the conditionality of transfers tied to children’s school attendance; other research has focused on the component of Mexico’s conditional cash transfer linked to health clinic attendance for maternal and child healthcare, but few studies have examined healthcare effects among older adults. We are aware of only two studies of health effects in CCT programs among the aging population. One study found that, two years after implementation, the Mexican CCT led to an average 4\% reduction in municipal-level mortality rates for those over age 65, primarily due to decreases in infectious diseases but also in diabetes-related death \citep{BarhamRowberry2013JDE}. The other study suggested that health utilization outcomes improved for female elderly, but not the male elderly six years after program implementation \citep{BehrmanParker2013Demography}. Previous research on non-contributory pensions or unconditional cash transfer programs (UCT) has also shown similarly mixed results. In Mexico, the UCT led to improved mental health 12 months after introduction \citep{GalianiGertlerBando2016LE, SalinasRodriguez2014PONE}, but greater mortality from cardiovascular and circulatory diseases was explained by changes in diet, increased in body mass index, and hypertension among recipients. In Yucatan of Mexico, the UCT improved verbal recall and lung function and reduced anemia six months after introduction \citep{Aguila2015PNAS}. In South Africa, one study found that a UCT improved self-reported health \citep{Case2004PerspectivesAging}, while another found it increased hypertension awareness but did not affect self-reported health or quality of life \citep{LloydSherlockAgrawal2014JDS}. Research in the United States has found that the supplementary security income (SSI) program positively affected elderly recipients’ health in the short run \citep{taubman1983supplemental}, but long-term effects are not well understood.\footnote{A recent survey \citep{richterman2023effects} on the effect of cash transfers on child and elderly mortality notes as follows: “Adults over the age of 60 made up only 1\ of our adult dataset. Our findings may therefore not apply to these older age groups.” } 

Mexico has high poverty rates, particularly among the elderly. According to the OECD, 24.7\% of Mexican adults aged 66 or older had less than half the median household income, partly due to the lack of universal coverage of the social security system \citep{Angel2017Gerontologist}. Furthermore, in 2017, Mexico ranked highest worldwide in the proportion of the population that is overweight and second highest in obesity prevalence, just below the United States \citep{OECD2017ObesityUpdate}. Miranda and colleagues found that the proportion of Mexican women who are overweight doubled from 1988 to 2012 and for men increased nearly 10 percentage points from 2000 to 2012 \citep{MirandaParkerSindelar2016ASHEcon}. This is especially relevant given that the income transfers from the CCT (Progresa) may have unintentionally increased adult obesity \citep{FernaldGertlerHou2008JN, KronebuschDamon2019EHB}.

In this paper, we analyze the effects of Progresa on older adults overall and by sex. Income-support programs may impact mortality through a variety of mechanisms. These include improved financial access to healthcare, especially in a context where many recipients were uninsured in the early years when most program expansions occurred. Increased cash may also affect non-medical pathways such as food consumption, potentially improving health among undernourished older adults, but it could also increase obesogenic food intake in a context marked by rapidly increasing obesity. Indeed, some evidence suggested that Progresa participation was associated with adult obesity  \citep{FernaldHouGertler2008PCD}, while other evidence found that higher cumulative Progresa cash transfers led to worse adult obesity \citep{FernaldGertlerHou2008JN}. Other mechanisms that might improve mortality outcomes include the conditionalities that promote regular health checkups, leading to better health behaviors and lower mortality. The distribution of resources and bargaining power to women may result in a greater emphasis on using available resources for health and nutrition than other uses. Given this conflicting evidence on mechanisms, and the potential for short-run versus long-run effects, evidence on the time-path of mortality impacts is particularly valuable. 
	
We combine nearly twenty years of nationwide vital statistics on mortality and  administrative records on Progresa receipt at the municipality level to estimate difference-in-difference models of the impacts on mortality. INEGI, Mexico’s statistical bureau, maintains individual-level vital statistics with detailed information on place of death and residence at the time of death, as well as health insurance coverage at the time of death. These variables allow us to link administrative data on Progresa receipt to the municipality of residence of the deceased. 


\subsection{Mexico's \textit{Progresa} Conditional Cash Transfer Program}\label{s:progresa}

The key feature of Progresa, which began in small rural communities in 1997, is providing monetary transfers to poor families. Some transfers are conditional on children attending school, while others require family members to visit health clinics. Progresa quickly expanded to cover six million families, which is one-fourth of all families in Mexico. Although the program expanded into urban areas, it remained largely rural, with two-thirds of its household beneficiaries living in communities with fewer than 2,500 inhabitants. The average family in Progresa received about MXN \$800 (US \$90.20 PPP) monthly, which increased household income by more than 25\% for many poor households. The program is means-tested, with both geographic and household-level targeting and all grants are given to the female head of the family.

The original program included two main types of grants: those linked to school attendance of children under 22 years of age in the household, and those linked to health and nutrition. The healthcare and nutritional component provides basic health care for all family members, with an emphasis on preventive health care. These services are provided by Mexican public health institutions, including the Ministry of Health and the Mexican Social Security Institute. To receive the fixed health and nutrition transfer, all members of beneficiary families must adhere to a regular schedule of health clinic visits. The calendar of visits varies by the age and sex of each individual, with visits required once per year for adults aged 17 and older, including the aging. Designated beneficiaries, generally mothers, are also required to attend monthly talks at the clinics on topics such as nutrition, hygiene, infectious diseases, immunization, family planning, and the detection and prevention of chronic diseases.  In 2006, a pension for the elderly was added to the program, providing a monthly payment to each adult aged 70 or older who is part of a Progresa family. 

The majority of elderly individuals in rural households eligible for Progresa live in extended families \citep{BehrmanParker2013Demography}. Before the program, about 75\% of elderly individuals lived in households with children under the age of 15. In most cases, the elderly individuals are the grandparents of children receiving Progresa educational grants and thus would not directly control its expenditures but still may benefit.  In a minority of households where elderly individuals live on their own, the female elderly would likely be the direct recipient of the Progresa transfers. 

The effects of the program on health may differ depending on whether the recipient of the transfer is an elderly member or a daughter/daughter-in-law (or another female member). Households where the elderly individuals live with children eligible to receive educational subsidies are likely to receive higher total transfers from PROGRESA/Oportunidades because these households are eligible for both payments conditional on health clinic visits and payments conditional on school enrollment and attendance. Additionally, who receives the transfer (the elderly person or a younger female) may affect how the expenditure is distributed, impacting the health of the elderly \citep{BehrmanParker2013Demography}.

We study data up until 2006 because after that, a new component of Progresa was implemented, which provided pension transfers directly to the elderly, fundamentally changing the benefits of the program for the aging. We concentrate only on the program during its first 10 years of operation in order to retain comparability and to estimate the impacts of the conditional cash transfer program rather than the conditional cash transfer program plus the pension for the aging. 

\subsection{Recent Trends in Over 65 Mortality in Mexico}\label{s:progresa}

Figure 1 shows trends in mortality for adults aged 65 and over overall and by sex in non-marginalized (non-poor) and marginalized (poor) areas from 1990 to 2016. In both areas, elderly mortality has been decreasing over time and is slightly higher in non-marginalized areas than marginalized areas.  Differences over time in aggregate mortality between marginalized and non-marginalized areas might reflect selective migration, where less healthy individuals migrate to urban areas, for instance, in search of additional health care facilities.  However, \cite{BeltranSanchezWong2025CSP}, using panel data from the Mexican Health and Aging Survey (MHAS) show that mortality by 2015 for those living in rural areas in 2001 (<2500) is significantly less than those living in more urban areas in 2001, suggesting selective migration does not explain mortality being higher in less urban areas.\footnote{\cite{BarhamRowberry2013JDE} also showed that elderly mortality was lower in Progresa municipalities than for Mexico as a whole (using raw vital statistics for the period 1992-2002).} 
\section{Data and Methodology}\label{s:data}

 \subsection{Data Sources}\label{s:sources}

\textit{\textbf{Mortality Data}} We use individual-level death registries derived from death certificates. These data provide each individual's cause of death, birth year, and sex. It comprises every death in the country between 1992 and 2016, with over three million records. Furthermore, the records include the municipality where the deceased resided, which we use to link to population denominators and program variables. All data is publicly available from the Mexican National Institute of Statistics and Geography (INEGI). INEGI is the national statistical agency and responsible for population censuses, collection of vital statistics, and a host of national surveys.\footnote{It derives mortality data from a certification system provided by the Mexican Ministry of Public Health. Data quality is a crucial concern when using the vital statistics registry. and Mexico generally maintains high-quality vital statistics \citep{mahapatra2007civil, mikkelsen2015global}}

The expansion of Mexico’s public health insurance program, \textit{Seguro Popular}, which continued to grow throughout the roll-out of \textit{Progresa}, may have improved the accuracy of death reporting by reducing financial barriers. This could have led to more accurate and timely documentation of deaths and their causes.\footnote{\textit{Seguro Popular} was introduced in 2002 in select municipalities and expanded in a quasi-random manner across Mexico, reaching full municipal coverage by 2006. However, significant portions of the eligible population remained uninsured, prompting continued enrollment efforts in subsequent years.} To mitigate bias from changes in vital statistics we control for the \textit{Seguro Popular} eligibility of a specific locality given the municipality's rollout, which addresses improvements in death registration during the period of our study.



\textit{\textbf{Population and Household Data}} 
We use INEGI population counts disaggregated by sex, municipality, and age estimated from census data (1990, 2000, 2005, 2010, 2020), which are also publicly available in INEGI. These data are used to estimate population exposure and are combined with our deaths records at the municipality level for each age group. To construct annual estimates, we linearly interpolate population counts for each sex and age group cell at the municipality level between 1990 and 2006. Similar information was drawn and interpolated on the number of households at the municipality level.

\textit{\textbf{\textit{Progresa} Administrative Data}} 
To compute our treatment variable, we used the number of enrolled households in \textit{Progresa} by year and by municipality, provided by \textit{Progresa} administrative personnel.

\textit{\textbf{Income and Expenditure Data}} We use the Mexican National Survey of Household Income and Expenditure (ENIGH) from 1992 to 2006. The survey provides a comprehensive overview of household income and expenditure by amount, source, and category; as well as occupational and sociodemographic characteristics, and housing infrastructure and household equipment. Specifically, we use detailed information on individual trimestral income sources (such as wages, independent business earnings, and transfers) and household trimestral spending on food and health, transportation, education, and clothing. Likewise, it has detailed information on household savings and individual non-monetary sources of income. All surveys are publicly available from INEGI.


 \textit{\textbf{Additional Controls}} We also used a municipality-level poverty index, the Margination Index in 1990, constructed by the Mexican Population Council (CONAPO), and interpolated for the years when the data were not available. The Margination Index classified all Mexican municipalities into 5 categories ranging from very low poverty to very high poverty, based on variables such as proportion of households with dirt floors and individual socioeconomic indicators.\footnote{The Margination Index was based on the share living in communities with less than 5000 inhabitants, the share earning less than twice the minimum wage, the share illiterate, and the shares with less than primary school, without a toilet, without electricity, without running water, with crowding, and with a dirt floor.} Second, as previously mentioned, we incorporate administrative data on \textit{Seguro Popular} beneficiaries to include an intensity variable at the municipality level to address coverage by this low-income health insurance program.\footnote{\textit{Seguro Popular} began a staggered rollout in 2002 at the municipal level.} Finally, we include the locality-level marginality index from the Mexican National Population Council, which serves as a proxy for poverty and socioeconomic conditions. 



% To account for potential confounding factors, we control for other programs and health services that were in place during the study period and could influence our results. First, we use administrative records on health infrastructure from 2002 to 2011, which include the number of clinics, hospitals, mobile clinics, health brigades, medical residents, and frontline healthcare personnel (doctors and nurses) at the locality level. These data are publicly available and provided by the Ministry of Health (Secretaría de Salud, 2022). 

\subsection{Sample Construction and Descriptive Statistics}\label{s:sample}

Our sample is restricted to deaths among individuals aged 65 and older for the mortality data to capture the program's effects. Then, we construct cells representing the number of deaths at the municipality level for each age and year. We combine these data with the population and household counts and the aforementioned time-varying municipality controls. Our sample of analyses focused on marginalized municipalities (high and very high poverty) identified as eligible by Progresa in 1997.


% As previously mentioned, we exclude the 2008 and 2009 expansion localities because (i) the focus of our study is 2007 expansion localities in rural localities, the most economically vulnerable areas, and for which we hypothesize that any mortality effects would be largest; (ii) other state non-contributory pension programs also covered localities in the 2008 expansion; and (iii) the 2009 expansion, which covered localities with 20,000--30,000 inhabitants, had insufficient take-up to capture the full program effect and could bias our result downward.\footnote{According to SEDESOL, the program covered approximately 1.03 millions of beneficiaries in its first year of expansion, it grew 80\% in the following expansion in 2008, and 10\% in the 2009 expansion.} We also end the analysis in 2011, as the 2012 expansion covered all localities, including large urban localities, preventing a valid control group after 2012. Following \cite{Feeney2018, feeney2017cash}, we use localities between 30,000 and 100,000 inhabitants as our preferred control group.\footnote{The results are not sensitive to including a higher inhabitants cutoff point.} This restriction avoids comparisons to large urban centers where mortality patterns may differ substantially from smaller, less populated regions. Consequently, the analysis focuses on deaths in rural localities eligible in 2007.\footnote{Furthermore, we decided not to pool the 2007 and 2008 group of expansion altogether because the recent literature on two-way fixed effects estimators has shown that estimates from linear models can differ from the group's average treatment on the treated (ATT) in the presence of treatment effect heterogeneity \citep{de2020two,callaway2021difference,sun2021estimating}. Therefore, in the absence of recent developments for heterogeneous treatment effects in non-linear models, we focus all the analyses separately for the 2007 expansion.} 

% We present summary statistics of deaths in Mexico for older adults aged 60 to 79 years old between 2002 and 2011, both overall and by eligibility group based on population locality size (Appendix Table \ref{at:death_stats}). There are on average 0.5 deaths per locality-age group in localities with fewer than 2,500 inhabitants (column 1). In localities eligible in the first expansion (year 2007), there were a total of 263,619 deaths across the study period (column 3), which accounted for over 15\% of deaths nationally. One fourth of these deaths were attributed to CVD causes, as shown in column (2), which was the largest category among all causes. 

Table 1 reports summary statistics of marginalized and non-marginalized municipalities in Mexico pre-program as well as the proportion of municipalities with Progresa beneficiaries. As the table demonstrates, marginalized municipalities have much poorer living conditions and as expected, a much higher rate of enrollment in Progresa.  


 Following the mortality results, we present analyses from the ENIGH household survey data to explore mechanisms that could help illuminate the mortality findings. Specifically, we use the all survey waves available between 1992 and 2006 to examine income and expenditure behaviors, providing three waves before and seven after the start of the \textit{Progresa} program.\footnote{The survey is available every two years in the period of analysis, except for an additional survey in 2005.} From these, we create a harmonized repeated cross-sectional dataset at the individual level, selecting adults ages 65 and above. We define three main income outcomes: individual earnings from wages, total household income, and \textit{Progresa} transfer income. Expenditures in the ENIGH are only measured at the household level; thus, we define total household spending, food consumption, and health spending (drugs and medical services). We also define specific categories such as cereals (grain, corn, and starches), meat and dairy, and sugar-related food (pastries, sugars, soft drinks, etc.) to test whether the cash transfer affects consumption of obesogenic foods rich in processed carbohydrates, starches, sugars, and fats. Additionally, we define alcohol and tobacco consumption as these are major risk factors for CVD and CVD-related mortality. From the individual occupation and sociodemographic characteristics, we identify each household member's employment status, and weekly hours worked. All income and expenditure monetary outcomes are converted to monthly values, updated to 2025 Mexican pesos, and then converted to U.S. dollars for comparability.\footnote{All yearly consumer price index and exchange rates used are publicly available from the Bank of Mexico webpage.}


% We present summary statistics pre- and post-program for eligible and ineligible localities. First, we document the presence of the \textit{70 y Más} cash transfer, which averages approximately \$15 USD in eligible localities during the post-expansion program period, and verify virtually no program transfer income in the pre-period and ineligible localities (Appendix Table \ref{at:ind_stats}).
% Regarding labor market outcomes, we observe a potential decline in earnings and income in eligible localities following the program’s introduction, consistent with reductions in employment and hours worked. As expected, enrollment in \textit{Seguro Popular} increases in eligible localities, aligning with the program’s municipal rollout during the same period.
% In terms of household-level indicators, we do not observe a relevant change in total spending among eligible localities (Appendix Table \ref{at:hh_stats}); however, there is a suggestive increase in food expenditures--—particularly in categories associated with obesogenic foods, including meat and dairy, sugar and fat, and cereals (Appendix Table \ref{at:hh_stats}). Finally, health-related spending appears to decrease from \$22.81 to \$13.12 (Appendix Table \ref{at:hh_stats_2}).
% Taken together, these patterns suggest plausible behavioral mechanisms through which the cash transfer may be influencing outcomes, and they merit further investigation. 


\subsection{Methods}\label{subsec:empirical_strategy}

We use a difference in difference estimation strategy,  using variation in the proportion of households benefiting from Progresa across municipalities, before and after Progresa began.  Following \cite{parker2023conditional}, we estimate impacts by comparing beneficiaries incorporated in the early (first) phase of Progresa 1997-1999 with those incorporated in the later, (second) 2000-2005 phase. The earlier phase corresponds to municipalities that enrolled households more intensively under Zedillo (1997–9) versus the later phase corresponds to those that enrolled households more intensively under Fox (2000–5). After 2005, the pattern of enrollment in Progresa changed substantially, moving towards more urban areas. 

The key explanatory variable of interest is the municipal proportion of household beneficiaries of Progresa incorporated into the program in the first phase between 1997 and 1999 interacted with a dummy measuring before and after Progresa implementation. We control also for a similar interaction with the proportion incorporated in 2000-2005.  Our estimates, therefore, reflect the impacts of being offered the program earlier during the first phase versus during the second phase of the program.  Note, this strategy relies on mortality being a cumulative variable, e.g., additional time receiving Progresa should lead to higher impacts on mortality. We also examine the differential impacts of the program by gender. 

Because Progresa was preferentially targeted toward more disadvantaged areas, a difference-in-differences design may be confounded by differential underlying trends between more and less disadvantaged municipalities. To mitigate this concern, we first restrict the sample to municipalities in the highest categories of the Mexican government’s marginality classification, where program roll-out was most intensive. Second, our specifications allow outcome differences between pre- and post-program cohorts to vary flexibly with cumulative municipal enrollment at the end of the program’s second roll-out phase. As a consequence, identification requires parallel trends only between early and late roll-out municipalities, rather than between early roll-out municipalities and all others \citep{parker2023conditional}. Third, we estimate event-study models that trace dynamic effects and test for differential pre-trends, using data that permit annual comparisons for up to seven years prior to the introduction of Progresa.
We estimate the following equation:

\begin{equation}
Mortality_{m,t}
= \beta_{0}\,Post_{m,t}\times Progresa_{m,1999}
+ \beta_{1},Post_{m,t}\times Progresa_{m,2005}
+ \delta_{m} + \gamma_{t} + \chi_{m,t} + \varepsilon_{m,t}
\end{equation}

where Mortality is the elderly mortality rate (65+) in municipality m in time t. We included interaction terms of Progresa intensity in 1999 or 2005 with Post Progresa dummy (equals 1 for 1997-2006 and 0 for 1991-1996). Fixed effects for municipality ($\delta_{m}$) and year ($\gamma_{t}$) were included to control for time-invariant unobservable municipality-level confounders and aggregate time trends. Covariate ($\chi_{m,t}$) is the proportion of beneficiaries for Seguro Popular, a nationwide health insurance program for informal sector workers.  We provide results for two specifications, one with and one without $\chi_{m,t}$. We also estimate event study specifications. 



% We next describe our empirical strategy for the mortality analysis, present the main estimates of the program's effect on mortality, and discuss robustness checks. 

% The phased rollout of the program, geographic coverage, and age eligibility allow us to implement a triple difference (difference-in-difference-in-differences) design, following \cite{feeney2017cash,Feeney2018}. In particular, we can use the timing of coverage among different localities and age eligibility to study changes in age-locality level outcomes before and after coverage. Since many small localities record large number of zero deaths for certain locality-age groups, the analysis relies on total death counts as the primary outcome, and models the population denominator controls for each age-locality cell as a population offset in the Poisson model. Recent literature recommends \citep{chen2024logs, Mullahy_Norton} the use of Poisson regression to model non-negative outcomes such as deaths, thereby avoiding outcome transformations to implement linear models with a large number of zero death rates. As a result, the Poisson regression naturally accommodates both zero and positive outcomes, assuming a consistent data-generating process across observations.\footnote{This approach treats zero and positive counts within a unified framework, without assuming separate processes as in zero-inflated or hurdle models. Furthermore, we acknowledge that Poisson regression assumes equality between the mean and variance, and that violations of this assumption could motivate the use of a negative binomial model. Nonetheless, given our large sample size and the robustness of the Poisson pseudo-maximum likelihood estimator (PPMLE), we are confident that our estimates of the conditional mean remain consistent even under overdispersion. Moreover, we cluster standard errors at the treatment level to account for any residual variance misspecification.} By including population as an offset, it models rates rather than raw counts, effectively adjusting for exposure. Importantly, it does so without requiring ad-hoc outcome transformations—such as $log(1+Y)$ or the inverse hyperbolic sine—that are unit-dependent and can distort results, and prevent interpretable effects as percentage changes.\footnote{\cite{Feeney2018, feeney2017cash}, implemented outcome transformations following the common practice at the time, that are currently considered to bias results, and hard to understand effects, which are largely analyzed in \cite{chen2024logs, Mullahy_Norton}}


% Therefore, following \cite{menares2025}, we define the outcomes of interest as the annual number of deaths within a locality for a given age (e.g., deaths in locality \textit{A} among people aged 60-69 in a given year). We then fit Poisson models for counts using a log link, including age-locality level population as an offset, with a general specification given by:

% \begin{align}
% y_{alt}=\exp\left(\beta\cdot70yMas_{alt} + \alpha_{a}+\gamma_{l} + \delta_{t} + \eta_{l_ea}+\xi_{l_et} + \theta_{ta} + \Pi_{lt} + Log(Pop_{alt})\right)\epsilon_{alt}
% \label{eq:did},
% \end{align}

% \noindent where $y_{alt}$ is the count of deaths for age-group $a$ in locality $l$ in year $t$.\footnote{It is important to notice that we have two age groups under the subscript \textit{a}, those between 60 to 69 and 70 to 79, which represent the ineligible age groups, respectively.} $70ymas_{alt}$ is an indicator that equals one from the first time a locality with age above 70 is eligible and onward, i.e., the treatment is an absorbing state. $\alpha_{a}$, $\gamma_{l}$, and $\delta_{t}$ are age group, locality, and year fixed effects, respectively. $l_e$ is a group locality eligibility indicator; thus, $\eta_{l_ea}$ represents fixed effects that control for time-invariant unobservable variables specific to localities-ages eligible groups, $\xi_{l_et}$ are time-variant fixed effects that account for year-specific eligible locality groups shocks common across age-groups, and $\theta_{ta}$ are time-variant fixed effects that account for age-eligible shocks common across localities. We also include a vector of controls, $\Pi_{lt}$, to account for time-variant locality characteristics, such as \textit{Progresa} penetration, \textit{Seguro Popular} coverage, health infrastructure, marginality index, and proxies for data quality, such as the average lag of death registration and the share of deaths certified by a medical doctor. Finally, $\epsilon_{alt}$ is the error term clustered (multi-way) by age and locality, the level of treatment. We weight all our regressions using the locality's 2005 population size age 60 and over. In this model, identification of the \textit{causal} effect of the \textit{70 y Más} program is predicated upon the assumption that---conditional on locality, age group, aggregate year shocks, time-invariant locality-age-eligible group, and time-variant locality-eligible and age-eligible cell indicators---there are no unobserved factors that correlated with both the timing and eligibility of coverage and other determinants of health outcomes. 


% In this case, identification of the causal effect of the program $\hat{\theta}_{ATT\%}$ is predicated upon the assumption that in the absence of treatment, the \textit{percentage changes} in the mortality rate would have been the same 

% To assess the plausibility of this parallel relative trends assumption, we examine the dynamic effects of \textit{70 y Más} using event studies around the time when the locality-age group becomes eligible. Given our data availability and focus on the initial 2007 rural expansion, we use five years before and four years after treatment (i.e., $t=-5$ in 2002, and $t=4$ in 2011).\footnote{We cannot analyze data before 2002 because of the absence of locality identifiers in the vital statistics data available to us. We end the analysis in 2011 because in 2012 the program further expanded to cover all localities in the country.} To estimate this dynamic version of our difference-in-difference-in-differences specification, we use a leads-and-lags model in event time, with the 2007 expansion year set to zero. Specifically, we estimate the following equation:

% \begin{align}
% y_{alt} = \exp\left(\sum_{k=-5}^{-2} \beta_k \cdot D_{alt}^k 
% + \sum_{k=0}^{3} \beta_k \cdot D_{alt}^k 
% + \alpha_{a} + \gamma_{l} + \delta_{t} 
% + \eta_{la} + \xi_{lt} 
% + \theta_{ta} + \Pi_{lt} + Log(Pop_{alt})\right) \cdot \epsilon_{alt}, 
% \label{eq:EventStudy}
% \end{align}


% \noindent where $D_{alt}^k= 1[t = 70yMas_{al} + k]$, and $70yMas_{al}$ is the timing of coverage of a locality eligible-age group $al$. In other words, $D_{alt}^k$ is a dummy variable indicating that a locality with deaths above the age 70, $al$, was covered by the program $k$ years ago (or included $k$ years ahead, for negative values of $k$). We normalize the coefficients such that $\beta_{k = -1} = 0$, i.e, treatment is re-coded in event time relative to the year before each locality age group was eligible for \textit{70 y Más}. Therefore, the $\beta_k$ coefficients can be interpreted as the effect of \textit{70 y Más} coverage on the mortality rate $y_{alt}$ for each $k$ period relative to the year before the coverage of a locality-age group $al$ in the \textit{70 y Más} program. 
 

\section{Mortality Impact of the \textit{Progresa} CCT Program}\label{s:reform_impact}

\subsection{Results: \textit{Progresa} Effect on Older Adult Mortality}\label{s:main_results}

Given the program’s delivery of benefits to women, we examine the differential effects of Progresa on mortality for individuals aged 65 and older in marginalized areas, both overall and by sex. Columns (1), (2), and (3) in Table 2 present the program's effect for both sexes combined, male only, and female only, respectively, while controlling for year and municipality fixed effects. Table 2 reports difference-in-differences estimates of the effect of Progresa on elderly (ages 65+) mortality in marginalized municipalities in Mexico between 1991 and 2006. The table presents estimates separately for men, women, and both sexes combined, with results shown for unweighted and population-weighted specifications, and with and without adjustment for the intensity of Seguro Popular.

Panel A shows unweighted estimates without controlling for Seguro Popular. The coefficient on the interaction between the post-program indicator and Progresa intensity is small and statistically insignificant for all three groups. Panel B adds controls for Seguro Popular intensity; the estimated effect of Progresa remains statistically insignificant and similar in magnitude.

Panels C and D present population-weighted specifications, where weights correspond to the size of the elderly population in each municipality. While the weighted estimates suggest more negative point estimates for male mortality, particularly in the post-Progresa period, these estimates remain statistically insignificant and are not mirrored for women or for the pooled sample. Controlling for Seguro Popular in Panel D does not materially alter the estimated effect of Progresa on elderly mortality.

Figures 2 and 3 present event study specifications for the specifications corresponding to Panels A (unweighted) and B (weighted) in Table 3.  While relatively noisy, the event study specifications confirm no significant effects of Progresa on mortality of the aging.  Pre program coefficients while again noisy in general do not show statistically significant pre program differences, providing support for our empirical strategy. In summary, across all specifications, both the event studies and the regression analysis provide no evidence of significant effects of Progresa on aging mortality over its first 10 years.


%  \textit{\textbf{Sex-Specific Event Studies}} 

% \textit{\textbf{Cause-Specific Event Studies}}  

\textit{\textbf{Robustness}} 


Our overall findings on Progresa’s effects are much less optimistic than a previous study by \citep{BarhamRowberry2013JDE}, which suggested that Progresa reduced municipal-level mortality among adults aged 65 or older by 4\%, primarily due to decreases in infectious diseases and also in diabetes-related deaths within two years of program implementation. To explore further, we replicate the specification of \cite{BarhamRowberry2013JDE}:

\begin{equation}
Mortality_{m,t}
= \alpha_{0}+\,Progresa_{m,t-2}
+ \delta_{m} + \gamma_{t} + \chi_{m,t} + \varepsilon_{m,t}
\end{equation}

where Mortality is the elderly mortality rate (65+) in municipality $m$ in time $t$ and Progresa is the second lag of program intensity and carry out several modifications of their specification (Appendix Table 1). Row A reports B\cite{BarhamRowberry2013JDE}’s results, where a one percentage point increase in 2-year lagged program intensity reduces approximately 0.064 deaths per 1,000 older adults. We replicate the specification, obtaining similar (but not identical) results in row B, with a negative and significant reduction of 0.07 deaths per 1,000 older adults (B).  \cite{BarhamRowberry2013JDE}’s results were unweighted, when we include population weights, statistical and negative significance remains but the size of impacts is attenuated by about 50\%. Using 1 year lagged program intensity (row D)  or 3 year lagged program intensity (row E), instead of \cite{BarhamRowberry2013JDE}’s 2 year lagged program intensity, further attenuates results to about one third of the original results and the results also become less precise estimates.

 
\subsection{Mechanisms: The \textit{Progresa} Program Impact on Work, Income and Household Expenditure}\label{s:mechanisms}


For the first time in Mexico, the \textit{Progresa} program provided a fixed amount of income to household heads complying with nationwide.\footnote{As previously discussed in section \ref{s:progresa}, in 2006, a pension for the elderly was added to the program, providing a monthly payment to each adult aged 70 or older who is part of a Progresa family.} 
% Additionally, there were other programs implemented at the state level, but these did not necessarily target rural areas or the whole nation. We find no significant differences in the trend once we remove states with other non-contributory pension programs before \textit{70 y Más} started (Appendix Figure \ref{af:es_robust}).} 
 Following its implementation, we do not find detectable effects of the cash transfer on mortality among early older adult living in eligible municipalities compared to later beneficiaries. Following previous literature, we initially hypothesized that analyzing deaths by cause could help illuminate underlying mechanisms. However, this section further investigates potential mechanisms using household income and expenditure survey data to examine whether older adults living in households exposed to the program exhibit income effects and thus also identify the main spending channels through which behavior changes may have occurred. 


\section{Discussion}\label{s:discussion}

This study examines the medium-term effect of a conditional cash transfer program on elderly mortality in poor areas of Mexico. We find no evidence of significant negative effects of the Progresa program to reduce aging mortality during its first 10 years. This is perhaps not surprising. Most of the cash in conditional cash transfer programs are linked to children’s school enrollment, effectively providing a subsidy to investing in children.  The presence of aging members of the household does not increase the amount of transfers received under Progresa. Further, because most aging individuals in Mexican households live in extended households, they are not likely to be the direct recipients of the cash transfers.  Finally, while they are required to attend preventive health clinic visits, the requirement for aging individuals is only one clinic visit, per year.  For all of these reasons, large and immediate effects on adult or elderly health and mortality might not be expected, compared with say pension payments to the aging. Mortality furthermore likely reflects cumulative health processes, and the primary channels through which Progresa operated—income smoothing, nutrition, and preventive care—are more likely to affect intermediate health outcomes in the short run.

We conclude that \cite{BarhamRowberry2013JDE}’s results are sensitive, becoming smaller and/or insignificant with relatively minor changes in the specification such as weighting or including different lags of intensity. Combined with the results presented here, we think the case for large negative effects of conditional cash transfers on adult mortality is not strong. 



\newpage
%%%%%%%%%%%%%%%%%%%%%%%%%%%%%%%%%%%%%%%%%%%%%%%%%%%%%%%%%%%%%%%%%%%%%%%%%%%%%%%%%%%%%%%%%%%%%%%
% REFERENCES
%%%%%%%%%%%%%%%%%%%%%%%%%%%%%%%%%%%%%%%%%%%%%%%%%%%%%%%%%%%%%%%%%%%%%%%%%%%%%%%%%%%%%%%%%%%%%%%
\setcounter{secnumdepth}{0}
%\clearpage
\bigskip
%\clearpage
\singlespace
\renewcommand{\bibname}{References}
%\bibliographystyle{apalike}
{\linespread{1.0}\selectfont\bibliography{bibliography}}

%%%%%%%%%%%%%%%%%%%%%%%%%%%%%%%%%%%%%%%%%%%%%%%%%%%%%%%%%%%%%%%%%%%%%%%%%%%%%%%%%%%%%%%%%%%%%%%
% FIGURES
%%%%%%%%%%%%%%%%%%%%%%%%%%%%%%%%%%%%%%%%%%%%%%%%%%%%%%%%%%%%%%%%%%%%%%%%%%%%%%%%%%%%%%%%%%%%%%%

\newpage
%=======================================================================
% figures.tex — Pre-Trend Event Study Figures (Mechanisms, ENIGH)
% Script: 05_mechanisms_eventstudy_enigh.do
% Layout: 2 columns x 4 rows per panel (max 8 subfigures)
% Paths:  pretrend/{marg,all}/F_PT_{1998,2000}_{outcome}
% Requires in main preamble: \usepackage{graphicx,subfigure}
%=======================================================================

% Compact subfigure helper: \ptsfig{filepath}{subcaption}
\newcommand{\ptsfig}[2]{%
  \subfigure[#2]{\includegraphics[width=0.47\linewidth]{#1}}}


%-----------------------------------------------------------------------


\begin{figure}[H]
    \caption[]{The effects of interaction terms between Progresa in 1999 and years on elderly mortality +65 overall and by sex in marginalized municipalities, Mexico, 1991-2006. }
    \label{f:es_dd_sex_no_weight}

  \subfigure[All]
{\includegraphics[width=0.69\textwidth]{figures/main/Figure_2a_all.pdf}}
    
    \subfigure[Males]
     {\includegraphics[width=0.49\textwidth]{figures/main/Figure_2b_male.pdf}}
    \subfigure[Females]
     {\includegraphics[width=0.49\textwidth]{figures/main/Figure_2c_female.pdf}}
     
         \floatfoot{Notes: The data analysis is restricted to only marginalized municipalities, 1991-2006. The plots present the interaction terms between Progresa intensity in 1999 and years. Year and municipality fixed effects are included. Standard errors are clustered at the municipality level. 1996 is a reference year and the regression models are not weighted.
Source: Authors’ calculations using data on deaths and population from the Instituto Nacional de Estadistica y Geografia and Mexican Census data.
}
\end{figure}

\begin{figure}[H]
    \caption[]{The effects of interaction terms between Progresa in 1999 and years on elderly mortality +65 overall and by sex in marginalized municipalities, Mexico, 1991-2006. }
    \label{f:es_dd_sex_no_weight}

  \subfigure[All]
{\includegraphics[width=0.69\textwidth]{figures/main/Figure_3a_all.pdf}}
    
    \subfigure[Males]
     {\includegraphics[width=0.49\textwidth]{figures/main/Figure_3b_male.pdf}}
    \subfigure[Females]
     {\includegraphics[width=0.49\textwidth]{figures/main/Figure_3c_female.pdf}}
     
         \floatfoot{Notes: The data analysis is restricted to only marginalized municipalities, 1991-2006. The plots present the interaction terms between Progresa intensity in 1999 and years. Year and municipality fixed effects are included. Standard errors are clustered at the municipality level. 1996 is a reference year and the regression models are not weighted.
Source: Authors’ calculations using data on deaths and population from the Instituto Nacional de Estadistica y Geografia and Mexican Census data.
}
\end{figure}


\begin{figure}[htbp]\centering
\ptsfig{figures/pretrend/all/F_PT_1998_employed}{Employment}\hfill
\ptsfig{figures/pretrend/all/F_PT_1998_hrs_worked}{Hrs Worked}\\[1ex]
\ptsfig{figures/pretrend/all/F_PT_1998_hrs_worked_pos}{Hrs Worked ($>$0)}\hfill
\ptsfig{figures/pretrend/all/F_PT_1998_ind_earnings}{Earnings}\\[1ex]
\ptsfig{figures/pretrend/all/F_PT_1998_ind_income_tot}{Income}\hfill
\ptsfig{figures/pretrend/all/F_PT_1998_progresa_ind}{PROGRESA Transfers}
\caption{Even Study: Individual Labor Outcomes - Intensity 1998, All Municipalities}
\label{fig:pt_1998_all_t1}
\end{figure}

\begin{figure}[htbp]\centering
\ptsfig{figures/pretrend/all/F_PT_1998_hh_earnings}{Earnings}\hfill
\ptsfig{figures/pretrend/all/F_PT_1998_hh_income_tot}{Income}\\[1ex]
\ptsfig{figures/pretrend/all/F_PT_1998_hh_expenditure}{Expenditure}\hfill
\ptsfig{figures/pretrend/all/F_PT_1998_progresa_hh}{PROGRESA}\\[1ex]
\ptsfig{figures/pretrend/all/F_PT_1998_savings}{Savings}\hfill
\ptsfig{figures/pretrend/all/F_PT_1998_debt}{Debt}
\caption{Even Study: Household Labor Outcomes - Intensity 1998, All Municipalities}
\label{fig:pt_1998_all_t2}
\end{figure}

\begin{figure}[htbp]\centering
\ptsfig{figures/pretrend/all/F_PT_1998_food_exp}{Total Food}\hfill
\ptsfig{figures/pretrend/all/F_PT_1998_vegg_fruit}{Veggies \& Fruit}\\[1ex]
\ptsfig{figures/pretrend/all/F_PT_1998_cereals}{Cereals}\hfill
\ptsfig{figures/pretrend/all/F_PT_1998_meat_dairy}{Meat \& Dairy}\\[1ex]
\ptsfig{figures/pretrend/all/F_PT_1998_sugar_fat_drink}{Sugar \& Fat}\hfill
\ptsfig{figures/pretrend/all/F_PT_1998_vice}{Alcohol and Tobacco}
\caption{Event Study: Household Food Outcomes - Intensity 1998, All Municipalities}
\label{fig:pt_1998_all_t3}
\end{figure}

\begin{figure}[htbp]\centering
\ptsfig{figures/pretrend/all/F_PT_1998_health_exp}{Total Health}\hfill
\ptsfig{figures/pretrend/all/F_PT_1998_medical}{Medical Visits}\\[1ex]
\ptsfig{figures/pretrend/all/F_PT_1998_medical_inpatient}{Inpatient}\hfill
\ptsfig{figures/pretrend/all/F_PT_1998_medical_outpatient}{Outpatient}\\[1ex]
\ptsfig{figures/pretrend/all/F_PT_1998_drugs}{Drugs (Rx and OTC)}\hfill
\ptsfig{figures/pretrend/all/F_PT_1998_ortho}{Orthotics}
\caption{Event Study: Household Health Outcomes - Intensity 1998, All Municipalities.}
\label{fig:pt_1998_all_t4}
\end{figure}






\begin{figure}[H]
    \caption[]{Share of progresa penetration}
    \label{f:es_dd_sex_no_weight}

    \subfigure[1997 - All]
     {\includegraphics[width=0.49\textwidth]{figures/appendix/FA0a_inten1997_kdensity.pdf}}
    \subfigure[1997 - by wave]
     {\includegraphics[width=0.49\textwidth]{figures/appendix/FA0b_inten1997_bywave.pdf}}
    

    \subfigure[1999 - All]
     {\includegraphics[width=0.49\textwidth]{figures/appendix/FA0c_inten1999_kdensity.pdf}}
    \subfigure[1999 - by wave]
     {\includegraphics[width=0.49\textwidth]{figures/appendix/FA0d_inten1999_bywave.pdf}}

    \subfigure[Continous - by Wave ENIGH Sample]
     {\includegraphics[width=0.49\textwidth]{figures/appendix/FA0f_intensity_new_surveytrend}}
    \subfigure[Continious - Mortality Sample]
     {\includegraphics[width=0.49\textwidth]{figures/appendix/FA0h_intensity_new_annual_trend}}
     
         \floatfoot{Notes: }
\end{figure}

%%%%%%%%%%%%%%%%%%%%%%%%%%%%%%%%%%%%%%%%%%%%%%%%%%%%%%%%%%%%%%%%%%%%%%%%%%%%%%%%%%%%%%%%%%%%%%%
% TABLES
%%%%%%%%%%%%%%%%%%%%%%%%%%%%%%%%%%%%%%%%%%%%%%%%%%%%%%%%%%%%%%%%%%%%%%%%%%%%%%%%%%%%%%%%%%%%%%%
\newpage


  

\begin{table}[H]
\RawCaption%
  {\caption[]{Summary statistics of non-marginalized and marginalized municipalities, Mexico (\%)}
   \label{t:descriptives}}
\floatbox[{\captop}]{table}[\FBwidth]
  {}
  {\scalebox{0.9}{
      \begin{tabular}{lcc}
\hline
 & \textbf{Non-marginalized areas} & \textbf{Marginalized areas} \\
\hline

\multicolumn{3}{l}{\textbf{A. Progresa program intensity}} \\
\quad 1999 & 0.10 (0.13) & 0.44 (0.22) \\
\quad 2005 & 0.28 (0.19) & 0.72 (0.17) \\[6pt]

\multicolumn{3}{l}{\textbf{B. Characteristics of municipalities, 1990}} \\
\multicolumn{3}{l}{\quad \textit{Components of the marginality index}} \\
\quad Illiterate population & 13.33 (6.57) & 33.33 (13.3) \\
\quad With less than primary education & 45.74 (12.32) & 69.56 (9.68) \\
\quad Without a toilet & 25.94 (16.62) & 60.30 (18.54) \\
\quad Without electricity & 11.81 (10.29) & 36.39 (24.72) \\
\quad Without running water & 19.43 (14.87) & 50.65 (24.07) \\
\quad With crowding & 59.84 (9.79) & 73.96 (8.24) \\
\quad With dirt floor & 22.05 (14.21) & 61.98 (21.58) \\
\quad In communities with $<$5{,}000 inhabitants & 60.18 (36.15) & 95.50 (13.54) \\
\quad Earning $<$2× minimum wage & 69.19 (11.60) & 85.82 (8.55) \\[6pt]

Number of municipalities & 1,234 & 1,119 \\
\hline
\end{tabular}

    }
    \vspace{-3mm}
    \floatfoot{Notes: Section A reports the mean proportion (\%) with standard deviations of municipalities with a ratio of the Progresa beneficiary households to the total number of households in the municipality. Section B presents sample mean (\%) of each variable and standard deviations are in parentheses. Marginalized areas include municipalities categorized as high and very high level of the Margination Index, while non-marginalized areas include municipalities categorized as medium, low, and very low level of the Margination Index. “With crowding” is measured by number of rooms divided by household size. Source: Authors’ calculations using Mexican Census 1990 and Progresa administrative data on beneficiaries.}
  }
\end{table}


\begin{table}[H]
\RawCaption%
  {\caption[]{Effects of Progresa on Elderly Mortality (65+) by Sex, Marginalized Municipalities}
   \label{t:did}}
\floatbox[{\captop}]{table}[\FBwidth]
  {}
  {\scalebox{0.9}{
      \begin{tabular}{lcccc} \hline \hline
& \multicolumn{1}{c}{(1)} & \multicolumn{1}{c}{(2)} & \multicolumn{1}{c}{(3)} & \multicolumn{1}{c}{(4)} \\ \toprule
\underline{\textit{Panel A: Pooled}}  \\ 
\textit{Intensity 1999 x post (1997-2006)} &        0.369 &        0.354 &       -0.777 &       -0.719 \\ 
 & (       2.546) & (       2.547) & (       1.464) & (       1.464) \\ 
  & & & & \\ 
\textit{Intensity 2005 x post (1997-2006)} &       -5.525* &       -5.513* &       -2.787 &       -2.796 \\ 
 & (       3.075) & (       3.074) & (       1.806) & (       1.803) \\ 
  & & & & \\ 
Mean (1991-1996) &        46.47 &        46.47 &        46.47 &        46.47 \\ 
Obs &       17,904 &       17,904 &       17,904 &       17,904 \\ 
  & & & & \\ 
\underline{\textit{Panel B: Males}}  \\ 
\textit{Intensity 1999 x post (1997-2006)} &       -1.627 &       -1.588 &       -2.288 &       -2.206 \\ 
 & (       2.767) & (       2.776) & (       1.732) & (       1.735) \\ 
  & & & & \\ 
\textit{Intensity 2005 x post (1997-2006)} &       -1.488 &       -1.519 &        0.214 &        0.201 \\ 
 & (       3.309) & (       3.312) & (       2.000) & (       1.994) \\ 
  & & & & \\ 
Mean (1991-1996) &        47.26 &        47.26 &        47.26 &        47.26 \\ 
Obs &       17,904 &       17,904 &       17,904 &       17,904 \\ 
  & & & & \\ 
\underline{\textit{Panel C: Females}}  \\ 
\textit{Intensity 1999 x post (1997-2006)} &        1.737 &        1.668 &        0.592 &        0.620 \\ 
 & (       2.932) & (       2.928) & (       1.824) & (       1.824) \\ 
  & & & & \\ 
\textit{Intensity 2005 x post (1997-2006)} &       -9.117** &       -9.063** &       -5.637** &       -5.641** \\ 
 & (       3.776) & (       3.770) & (       2.205) & (       2.205) \\ 
  & & & & \\ 
Mean (1991-1996) &        46.11 &        46.11 &        46.11 &        46.11 \\ 
Obs &       17,904 &       17,904 &       17,904 &       17,904 \\ 
  & & & & \\ 
No. Mun &        1,119 &        1,119 &        1,119 &        1,119 \\ 
  & & & & \\ 
Year FE & Y & Y & Y & Y \\ 
Mun FE & Y & Y & Y & Y \\ 
Seguro Popular & N & Y & N & Y \\ 
Weights & N & N & Y & Y \\ 
Cluster SE: Mun & Y & Y & Y & Y \\ 
\bottomrule
\end{tabular}
    }
    \vspace{-3mm}
    \floatfoot{Notes: The analysis is restricted to marginalized municipalities, 1991–2006. Progresa and Seguro Popular intensities are continuous variables. The post-dummy equals 1 for years 1997–2006. Standard errors are clustered at the municipality level and shown in parentheses. ***p$<0.01$, **p$<0.05$, *p$<0.1$.
}
  }
\end{table}


\begin{table}[htbp]
\centering
\caption{Progresa Impacts on Labor Supply and Living Arrangements: Eligible Elderly (Age 65+ in 1997)}
\label{tab:combined_elderly}
\begin{threeparttable}

\begin{tabular}{lcccc}
\toprule
 & (1) & (2) & (3) & (4) \\
 & Weekly Hours & Live Alone & With Children & Only Elderly \\
\midrule
\multicolumn{5}{l}{\textbf{Panel A: Females}} \\
\midrule
Treat $\times$ 1998 & 0.124 & 0.000 & 0.001 & $-$0.008 \\
                    & (1.136) & (0.008) & (0.011) & (0.008) \\[4pt]
Treat $\times$ 1999 & 0.292 & 0.007 & 0.015 & 0.012 \\
                    & (1.070) & (0.010) & (0.018) & (0.014) \\[4pt]
Observations        & 3,338 & 3,597 & 3,597 & 3,597 \\
Control Mean (1997) & 3.923 & 0.044 & 0.708 & 0.145 \\
\midrule
\multicolumn{5}{l}{\textbf{Panel B: Males}} \\
\midrule
Treat $\times$ 1998 & $-$1.210 & 0.012 & $-$0.008 & $-$0.007 \\
                    & (1.895) & (0.008) & (0.011) & (0.010) \\[4pt]
Treat $\times$ 1999 & 0.378 & 0.008 & 0.040$^{**}$ & $-$0.001 \\
                    & (1.872) & (0.012) & (0.020) & (0.017) \\[4pt]
Observations        & 3,170 & 3,419 & 3,419 & 3,419 \\
Control Mean (1997) & 25.288 & 0.049 & 0.676 & 0.159 \\
\bottomrule
\end{tabular}

\begin{tablenotes}
\small
\item \textit{Notes:} Difference-in-differences estimates with 1997 as baseline. Sample restricted to 
eligible households. Columns (1) reports weekly hours worked. Columns (2)--(4) report living 
arrangement outcomes: living alone, living with children, and living only with other elderly, 
respectively. Standard errors in parentheses.
\item $^{*}$ $p < 0.10$, $^{**}$ $p < 0.05$, $^{***}$ $p < 0.01$
\end{tablenotes}
\end{threeparttable}

\end{table}




\appendix
% \begin{bibunit}
%%%%%%%%%%%%%%%%%%%%%%%%%%%%%%%%%%%%%%%%%%%%%%%%%%%%%%%%%%%%%%%%%%%%%%%%%%%%%%%%%%%%%%%%%%%%%%%
% APPENDIX
%%%%%%%%%%%%%%%%%%%%%%%%%%%%%%%%%%%%%%%%%%%%%%%%%%%%%%%%%%%%%%%%%%%%%%%%%%%%%%%%%%%%%%%%%%%%%%%

\newpage
\setcounter{page}{1}
\section*{Online Appendix}\label{appOnline}
\vspace{2cm}
\begin{center}
{\Large \it Do CCT Programs (Really) Reduce Mortality: 10 Year Evidence From Mexico} \\[.4cm]
Emma Aguila, William H. Dow, Felipe Menares, Susan Parker, Jorge Peniche, and Soomin Ryu
\end{center}
\listoffigures
\listoftables
%%%%%% APPENDIX A
\newpage
\section*{Appendix A: Additional Figures and Tables}\label{appA}
\setcounter{table}{0}
\renewcommand{\thetable}{A.\arabic{table}}
\setcounter{figure}{0}
\renewcommand{\thefigure}{A.\arabic{figure}}
\addtocontents{toc}{\setcounter{tocdepth}{1}}
\renewcommand{\thesubsection}{\Alph{subsection}}



\begin{figure}[H]
    \caption{Mortality and Progresa Intensity in Mexico}
    \label{f:variation}

    \subfigure[Mortality and Progresa Intensity]
    {\includegraphics[width=0.69\textwidth]{figures/appendix/Figure_1a_all.pdf}}

    \subfigure[Intensity in 1999 - All Municipalities]
     {\includegraphics[width=0.49\textwidth]{figures/appendix/Figure_1b_inten1999_all.pdf}}
    \subfigure[Intensity in 2005 - All Municipalities]
     {\includegraphics[width=0.49\textwidth]{figures/appendix/Figure_1c_inten2005_all.pdf}}
      \subfigure[Intensity in 1997 - Highly Marginalized Municipalities]
     {\includegraphics[width=0.49\textwidth]{figures/appendix/Figure_1d_inten1997_mort.pdf}}
    \subfigure[Intensity in 1997 - All Municipalities]
     {\includegraphics[width=0.49\textwidth]{figures/appendix/Figure_1e_inten1997_mort_all.pdf}}
    
     
         \floatfoot{Notes: }
\end{figure}

\begin{figure}[H]
    \caption{The effects of interaction terms between Progresa in 1999 and years on elderly mortality +65 overall and by sex in marginalized municipalities, Mexico, 1991-2006. }
    \label{f:es_dd_sex_no_weight}

  \subfigure[Unweighted]
{\includegraphics[width=0.89\textwidth]{figures/appendix/Figure_2a_uw.pdf}}
  \subfigure[AAMR Weighted +SP]
{\includegraphics[width=0.89\textwidth]{figures/appendix/Figure_3_aamr_col4.pdf}}
    
  

    
         
         \floatfoot{Notes: The data analysis is restricted to only marginalized municipalities, 1991-2006. The plots present the interaction terms between Progresa intensity in 1999 and years. Year and municipality fixed effects are included. Standard errors are clustered at the municipality level. 1996 is a reference year and the regression models are not weighted.
Source: Authors’ calculations using data on deaths and population from the Instituto Nacional de Estadistica y Geografia and Mexican Census data.
}
\end{figure}



\begin{figure}[H]
    \caption{The effects of interaction terms between Progresa in 1999 and years on elderly mortality +65 overall and by sex in marginalized municipalities, Mexico, 1991-2006. }
    \label{f:es_dd_sex_no_weight}

  \subfigure[BR - EMR - All]
{\includegraphics[width=0.49\textwidth]{figures/appendix/Figure_3_emr65_BR.pdf}}
    \subfigure[HM - EMR - All]
{\includegraphics[width=0.49\textwidth]{figures/appendix/Figure_3_emr65_HighMarg.pdf}}


\subfigure[BR - EMR - Females]
{\includegraphics[width=0.49\textwidth]{figures/appendix/Figure_3_emr65f_BR.pdf}}
\subfigure[HM - EMR - Females]
     {\includegraphics[width=0.49\textwidth]{figures/appendix/Figure_3_emr65f_HighMarg.pdf}}

\subfigure[BR - EMR - Males]
{\includegraphics[width=0.49\textwidth]{figures/appendix/Figure_3_emr65m_BR.pdf}}
 \subfigure[HM - EMR - Males]
{\includegraphics[width=0.49\textwidth]{figures/appendix/Figure_3_emr65m_HighMarg.pdf}}
    
         \floatfoot{Notes: The data analysis is restricted to only marginalized municipalities, 1991-2006. The plots present the interaction terms between Progresa intensity in 1999 and years. Year and municipality fixed effects are included. Standard errors are clustered at the municipality level. 1996 is a reference year and the regression models are not weighted.
Source: Authors’ calculations using data on deaths and population from the Instituto Nacional de Estadistica y Geografia and Mexican Census data.
}
\end{figure}



% \begin{figure}[H]
%     \caption{The effects of interaction terms between Progresa in 1999 and years on elderly mortality +65 overall and by sex in marginalized municipalities, Mexico, 1991-2006. }
%     \label{f:es_dd_sex_no_weight}

%   \subfigure[BR - AAMR - All]
% {\includegraphics[width=0.49\textwidth]{figures/appendix/Figure_3_aamr65_BR.pdf}}
%     \subfigure[HM - AAMR - All]
% {\includegraphics[width=0.49\textwidth]{figures/appendix/Figure_3_aamr65_HighMarg.pdf}}


% \subfigure[BR - AAMR - Females]
% {\includegraphics[width=0.49\textwidth]{figures/appendix/Figure_3_aamr65f_BR.pdf}}
% \subfigure[HM - AAMR - Females]
%      {\includegraphics[width=0.49\textwidth]{figures/appendix/Figure_3_aamr65f_HighMarg.pdf}}

% \subfigure[BR - AAMR - Males]
% {\includegraphics[width=0.49\textwidth]{figures/appendix/Figure_3_aamr65m_BR.pdf}}
%  \subfigure[HM - AAMR - Males]
% {\includegraphics[width=0.49\textwidth]{figures/appendix/Figure_3_aamr65m_HighMarg.pdf}}
    
%          \floatfoot{Notes: The data analysis is restricted to only marginalized municipalities, 1991-2006. The plots present the interaction terms between Progresa intensity in 1999 and years. Year and municipality fixed effects are included. Standard errors are clustered at the municipality level. 1996 is a reference year and the regression models are not weighted.
% Source: Authors’ calculations using data on deaths and population from the Instituto Nacional de Estadistica y Geografia and Mexican Census data.
% }
% \end{figure}




\begin{figure}[H]
    \caption{The effects of interaction terms between Progresa in 1999 and years on elderly mortality +65 overall and by sex in marginalized municipalities, Mexico, 1991-2006. }
    \label{f:es_dd_sex_no_weight}

  \subfigure[Cancer]
{\includegraphics[width=0.49\textwidth]{figures/appendix/Figure_4_tb_cancer_Marg.pdf}}
\subfigure[Cardiovascular]
{\includegraphics[width=0.49\textwidth]{figures/appendix/Figure_4_tb_card_Marg.pdf}}

    \subfigure[Respiratory]
{\includegraphics[width=0.49\textwidth]{figures/appendix/Figure_4_tb_resp_Marg.pdf}}
    \subfigure[Infectious]
     {\includegraphics[width=0.49\textwidth]{figures/appendix/Figure_4_tb_infect_Marg.pdf}}
    

     
         \floatfoot{Notes: The data analysis is restricted to only marginalized municipalities, 1991-2006. The plots present the interaction terms between Progresa intensity in 1999 and years. Year and municipality fixed effects are included. Standard errors are clustered at the municipality level. 1996 is a reference year and the regression models are not weighted.
Source: Authors’ calculations using data on deaths and population from the Instituto Nacional de Estadistica y Geografia and Mexican Census data.
}
\end{figure}


\begin{figure}[H]
    \caption{The effects of interaction terms between Progresa in 1999 and years on elderly mortality +65 overall and by sex in marginalized municipalities, Mexico, 1991-2006. }
    \label{f:es_dd_sex_no_weight}



         \subfigure[Diabetes]
{\includegraphics[width=0.49\textwidth]{figures/appendix/Figure_4_tb_diab_Marg.pdf}}
    \subfigure[Nutri]
     {\includegraphics[width=0.49\textwidth]{figures/appendix/Figure_4_tb_nutri_Marg.pdf}}

              \subfigure[Ill-Defined]
{\includegraphics[width=0.49\textwidth]{figures/appendix/Figure_4_tb_illdef_Marg.pdf}}
    \subfigure[Accidents]
     {\includegraphics[width=0.49\textwidth]{figures/appendix/Figure_4_tb_accid_Marg.pdf}}

     

     
         \floatfoot{Notes: The data analysis is restricted to only marginalized municipalities, 1991-2006. The plots present the interaction terms between Progresa intensity in 1999 and years. Year and municipality fixed effects are included. Standard errors are clustered at the municipality level. 1996 is a reference year and the regression models are not weighted.
Source: Authors’ calculations using data on deaths and population from the Instituto Nacional de Estadistica y Geografia and Mexican Census data.
}
\end{figure}



\begin{figure}[H]
    \caption{The effects of interaction terms between Progresa in 1999 and years on elderly mortality +65 overall and by sex in marginalized municipalities, Mexico, 1991-2006. }
    \label{f:es_dd_sex_no_weight}

  \subfigure[Cancer]
{\includegraphics[width=0.49\textwidth]{figures/appendix/Figure_4_tb_cancer_BR.pdf}}
\subfigure[Cardiovascular]
{\includegraphics[width=0.49\textwidth]{figures/appendix/Figure_4_tb_card_BR.pdf}}

    \subfigure[Respiratory]
{\includegraphics[width=0.49\textwidth]{figures/appendix/Figure_4_tb_resp_BR.pdf}}
    \subfigure[Infectious]
     {\includegraphics[width=0.49\textwidth]{figures/appendix/Figure_4_tb_infect_BR.pdf}}


     

     
         \floatfoot{Notes: The data analysis is restricted to only marginalized municipalities, 1991-2006. The plots present the interaction terms between Progresa intensity in 1999 and years. Year and municipality fixed effects are included. Standard errors are clustered at the municipality level. 1996 is a reference year and the regression models are not weighted.
Source: Authors’ calculations using data on deaths and population from the Instituto Nacional de Estadistica y Geografia and Mexican Census data.
}
\end{figure}


\begin{figure}[H]
    \caption{The effects of interaction terms between Progresa in 1999 and years on elderly mortality +65 overall and by sex in marginalized municipalities, Mexico, 1991-2006. }
    \label{f:es_dd_sex_no_weight}

         \subfigure[Diabetes]
{\includegraphics[width=0.49\textwidth]{figures/appendix/Figure_4_tb_diab_BR.pdf}}
    \subfigure[Nutri]
     {\includegraphics[width=0.49\textwidth]{figures/appendix/Figure_4_tb_nutri_BR.pdf}}

              \subfigure[Ill-Defined]
{\includegraphics[width=0.49\textwidth]{figures/appendix/Figure_4_tb_illdef_BR.pdf}}
    \subfigure[Accidents]
     {\includegraphics[width=0.49\textwidth]{figures/appendix/Figure_4_tb_accid_BR.pdf}}

     

     
         \floatfoot{Notes: The data analysis is restricted to only marginalized municipalities, 1991-2006. The plots present the interaction terms between Progresa intensity in 1999 and years. Year and municipality fixed effects are included. Standard errors are clustered at the municipality level. 1996 is a reference year and the regression models are not weighted.
Source: Authors’ calculations using data on deaths and population from the Instituto Nacional de Estadistica y Geografia and Mexican Census data.
}
\end{figure}




%%%%%%%%%%%%%%%%%%%%%%%%%%%%
%%%%%%%%%%%TABLES
%%%%%%%%%%%%%%%%%%%%%%%%%%%%



\newcommand{\mectable}[4]{%
  \begin{table}[H]
    \centering
    \caption{#1}
    \label{#2}
    \resizebox{\textwidth}{!}{\input{#3}}
    {\footnotesize\raggedright\noindent Notes: #4\par}
  \end{table}
  \clearpage
}

\begin{table}[H]
\RawCaption%
  {\caption{\deleted{Barham and Rowberry’s results and our replication: Impact of Progresa on elderly mortality 65+, municipalities incorporated in 1998 or 1999}\added{Replication of Barham and Rowberry (2013): Effect of Progresa on Elder Mortality Rate (EMR 65+), Municipalities Incorporated in 1998--1999}}
   \label{t:did_robust_br}}
\floatbox[{\captop}]{table}[\FBwidth]
  {}
  {\scalebox{0.9}{
      \begin{tabular}{lccc} \hline \hline
& \multicolumn{1}{c}{Pooled} & \multicolumn{1}{c}{Males} & \multicolumn{1}{c}{Females} \\ 
\cmidrule(lr){2-2}\cmidrule(lr){3-3}\cmidrule(lr){4-4}  
& \multicolumn{1}{c}{(1)} & \multicolumn{1}{c}{(2)} & \multicolumn{1}{c}{(3)} \\ \toprule 
\underline{\textit{Panel A: BR (2013)}}  \\ 
\textit{2-yr lagged Intensity} & -6.37*** & -6.42*** & -6.46*** \\ 
 & (1.04) & (1.42) & (1.31) \\ 
  & & & \\ 
Mean 1996 & 47.5 & 49.3 & 46.0 \\ 
Obs & 21,571 & 21,571 & 21,571 \\ 
No. Mun & 1,961 & 1,961 & 1,961 \\ 
  & & & \\ 
\underline{\textit{Panel B: Replication (Unweighted)}}  \\ 
\textit{2-yr lagged Intensity} &       -7.061*** &       -6.889*** &       -7.426*** \\ 
 & (       1.256) & (       1.711) & (       1.559) \\ 
  & & & \\ 
\underline{\textit{Panel C: Replication (Weighted)}}  \\ 
\textit{2-yr lagged Intensity} &       -3.740*** &       -3.322*** &       -4.210*** \\ 
 & (       0.786) & (       0.935) & (       1.002) \\ 
  & & & \\ 
\underline{\textit{Panel D: 1-yr Lag (Weighted)}}  \\ 
\textit{1-yr lagged Intensity} &       -3.457*** &       -2.681** &       -4.255*** \\ 
 & (       0.890) & (       1.072) & (       1.043) \\ 
  & & & \\ 
\underline{\textit{Panel E: 3-yr Lag (Weighted)}}  \\ 
\textit{3-yr lagged Intensity} &       -2.071*** &       -2.357** &       -1.830* \\ 
 & (       0.772) & (       0.957) & (       0.960) \\ 
  & & & \\ 
Mean 1996 &        48.16 &        48.87 &        47.49 \\ 
Obs &       15,642 &       15,642 &       15,642 \\ 
No. Mun &        1,422 &        1,422 &        1,422 \\ 
  & & & \\ 
Year FE & Y & Y & Y \\ 
Mun FE & Y & Y & Y \\ 
Cluster SE: Mun & Y & Y & Y \\ 
\bottomrule
\end{tabular}
    }
    \vspace{-3mm}
    \floatfoot{\deleted{Notes: The data analysis is restricted to municipalities incorporated into Progresa in 1998 or 1999. Progresa intensity data are pooled from 1990 to 2000. All regression models include year fixed effects and municipality fixed effects. Standard errors are clustered at the municipality level. The regression model is weighted by the number of population aged +65 for both sexes in (1), only male in (2), and only female in (3).\\ ***p$<0.01$, **p$<0.05$, *p$<0.1$.}\added{Notes: This table replicates the main specification of Barham and Rowberry (2013). The sample is restricted to municipalities incorporated into Progresa in 1998 or 1999. The outcome is the Elder Mortality Rate (EMR) for adults aged 65 and older per 1,000 inhabitants. Progresa intensity is measured using data pooled from 1990 to 2000. All regressions include municipality and year fixed effects. The regression model is weighted by the population aged 65 and older, shown separately for the pooled (column 1), male (column 2), and female (column 3) populations. Standard errors are clustered at the municipality level and shown in parentheses. ***p$<0.01$, **p$<0.05$, *p$<0.1$.}}
  }
\end{table}
\newpage




\begin{table}[H]
\RawCaption%
  {\caption{\deleted{Barham and Rowberry Extended Results Replication}\added{Sensitivity Analysis: Barham and Rowberry (2013) Replication with Alternative Lag Structures and Sample Restrictions}}
   \label{t:did_robust_br2}}
\floatbox[{\captop}]{table}[\FBwidth]
  {}
  {\scalebox{0.9}{
      \input{tables/appendix/T_BR_robustness_emr65}
    }
    \vspace{-3mm}
    \floatfoot{\deleted{Notes: Progresa intensity data are pooled from 1990 to 2000. All regression models include year fixed effects and municipality fixed effects. Standard errors are clustered at the municipality level. The regression model is weighted by the number of population aged +65. ***p$<0.01$, **p$<0.05$, *p$<0.1$.}\added{Notes: This table reports sensitivity analyses extending the replication of Barham and Rowberry (2013), varying the lag structure of Progresa intensity, weighting, and sample restrictions. The outcome is the Elder Mortality Rate (EMR) for adults aged 65 and older per 1,000 inhabitants. Row (A) reproduces Barham and Rowberry's (2013) original results; row (B) reports our replication using the same specification. Subsequent rows vary the lag structure (1-year and 3-year lags), weighting, and sample. Progresa intensity data are pooled from 1990 to 2000. All regressions include municipality and year fixed effects. Standard errors are clustered at the municipality level. ***p$<0.01$, **p$<0.05$, *p$<0.1$.}}
  }
\end{table}

% \begin{table}[H]
% \RawCaption%
%   {\caption{Barham and Rowberry Extended Results Replication}
%    \label{t:did_robust}}
% \floatbox[{\captop}]{table}[\FBwidth]
%   {}
%   {\scalebox{0.9}{
%       \input{tables/appendix/T_BR_robustness_aamr65}
%     }
%     \vspace{-3mm}
%     \floatfoot{Notes: Progresa intensity data are pooled from 1990 to 2000. All regression models include year fixed effects and municipality fixed effects. Standard errors are clustered at the municipality level. The regression model is weighted by the number of population aged +65. ***p$<0.01$, **p$<0.05$, *p$<0.1$.}
%   }
% \end{table}





% \begin{table}[H]
% \RawCaption%
%   {\caption{Functional Forms}
%    \label{t:did_robust}}
% \floatbox[{\captop}]{table}[\FBwidth]
%   {}
%   {\scalebox{0.9}{
%       \begin{tabular}{lccc} \hline \hline
& Levels & Log & Poisson \\ 
\cmidrule(lr){2-2}\cmidrule(lr){3-3}\cmidrule(lr){4-4}
& \multicolumn{1}{c}{(1)} & \multicolumn{1}{c}{(2)} & \multicolumn{1}{c}{(3)} \\ \toprule
\underline{\textit{Panel A: Pooled}}  \\ 
\textit{Intensity 1999 x post (1997-2006)} &       -0.777 &       -0.020 &       -0.031 \\ 
  & (       1.464) & (       0.042) & (       0.058) \\ 
   & & & \\ 
\textit{Intensity 2005 x post (1997-2006)} &       -2.787 &       -0.052 &       -0.128* \\ 
 & (       1.806) & (       0.052) & (       0.072) \\ 
  & & & \\ 
Mean (1991-1996)  &        46.47 &         3.75 &        23.86 \\ 
  & & &  \\ 
\underline{\textit{Panel B: Males}}  \\ 
\textit{Intensity 1999 x post (1997-2006)} &       -2.288 &       -0.029 &       -0.104* \\ 
  & (       1.732) & (       0.050) & (       0.062) \\ 
  & & & \\ 
\textit{Intensity 2005 x post (1997-2006)}  &        0.214 &       -0.020 &       -0.024 \\ 
 & (       2.000) & (       0.057) & (       0.076) \\ 
 & & & \\ 
Mean (1991-1996) &        47.26 &         3.78 &        12.25 \\ 
  & & &  \\ 
\underline{\textit{Panel C: Females}}  \\ 
\textit{Intensity 1999 x post (1997-2006)}  &        0.592 &        0.025 &        0.055 \\ 
  & (       1.824) & (       0.055) & (       0.074) \\ 
  & & & \\ 
\textit{Intensity 2005 x post (1997-2006)} &       -5.637** &       -0.173** &       -0.223*** \\ 
 & (       2.205) & (       0.067) & (       0.076) \\ 
   & & & \\ 
Mean (1991-1996)  &        46.11 &         3.73 &        11.61 \\ 
  & & &  \\ 
Obs &       17,904 &       17,223 &       17,904 \\ 
No. Mun &        1,119 &        1,119 &        1,119 \\ 
  & & & \\ 
Year FE & Y & Y & Y \\ 
Mun FE  & Y & Y & Y \\ 
Weights & Y & Y & Y \\ 
Cluster SE: Mun & Y & Y & Y \\ 
\bottomrule
\end{tabular}
%     }
%     \vspace{-3mm}
%     \floatfoot{Notes: Progresa intensity data are pooled from 1990 to 2000. All regression models include year fixed effects and municipality fixed effects. Standard errors are clustered at the municipality level. The regression model is weighted by the number of population aged +65. ***p$<0.01$, **p$<0.05$, *p$<0.1$.}
%   }
% \end{table}
% \newpage






\begin{table}[H]
\RawCaption%
  {\caption{\deleted{AAMR}\added{Effects of Progresa Enrollment Intensity on Age-Adjusted Mortality Rate (AAMR 65+) by Sex, Marginalized Municipalities, 1991--2006}}
   \label{t:did_robust_aamr}}
\floatbox[{\captop}]{table}[\FBwidth]
  {}
  {\scalebox{0.9}{
      \input{tables/appendix/AT4_aamr_mortality}
    }
    \vspace{-3mm}
    \floatfoot{\deleted{Notes: Progresa intensity data are pooled from 1990 to 2000. All regression models include year fixed effects and municipality fixed effects. Standard errors are clustered at the municipality level. The regression model is weighted by the number of population aged +65. ***p$<0.01$, **p$<0.05$, *p$<0.1$.}\added{Notes: This table reports difference-in-differences estimates of the effect of Progresa enrollment intensity on the Age-Adjusted Mortality Rate (AAMR) for adults aged 65 and older per 1,000 inhabitants, separately for the pooled, male, and female populations. Columns (1) and (2) report unweighted estimates without and with Seguro Popular intensity; columns (3) and (4) report estimates weighted by the population aged 65 and older without and with Seguro Popular intensity. The post-dummy equals 1 for years 1997--2006. The sample is restricted to highly marginalized municipalities ($\textit{gm\_mun\_1990} = 4$ or $5$). All regressions include municipality and year fixed effects. Standard errors are clustered at the municipality level. ***p$<0.01$, **p$<0.05$, *p$<0.1$.}}
  }
\end{table}

\begin{landscape}
\begin{table}[H]
\RawCaption%
  {\caption{\deleted{AAMR}\added{Effects of Progresa Enrollment Intensity on Cause-Specific Elder Mortality Rate (EMR 65+), Highly Marginalized Municipalities, 1991--2006}}
   \label{t:did_robust_cod}}
\floatbox[{\captop}]{table}[\FBwidth]
  {}
  {\scalebox{0.7}{
      \input{tables/appendix/AT5_cod_mortality}
    }
    \vspace{-3mm}
    \floatfoot{\deleted{Notes: Progresa intensity data are pooled from 1990 to 2000. All regression models include year fixed effects and municipality fixed effects. Standard errors are clustered at the municipality level. The regression model is weighted by the number of population aged +65. ***p$<0.01$, **p$<0.05$, *p$<0.1$.}\added{Notes: This table reports difference-in-differences estimates of the effect of Progresa enrollment intensity on cause-specific Elder Mortality Rate (EMR) for adults aged 65 and older per 1,000 inhabitants. Each column corresponds to a specific cause of death. The post-dummy equals 1 for years 1997--2006. The sample is restricted to highly marginalized municipalities ($\textit{gm\_mun\_1990} = 4$ or $5$). All regressions include municipality and year fixed effects. Standard errors are clustered at the municipality level. ***p$<0.01$, **p$<0.05$, *p$<0.1$.}}
  }
\end{table}
\end{landscape}

\begin{landscape}
\begin{table}[H]
\RawCaption%
  {\caption{\deleted{AAMR}\added{Effects of Progresa Enrollment Intensity on Cause-Specific Elder Mortality Rate (EMR 65+), Barham-Rowberry Sample, 1992--2002}}
   \label{t:did_robust_cod_br}}
\floatbox[{\captop}]{table}[\FBwidth]
  {}
  {\scalebox{0.9}{
      \begin{tabular}{lcccccccccc} \hline \hline
& Cancer & Diab. & IllDef & Resp. & Card. & Infect. & Nutri. & Accid. & Other \\ 
\cmidrule(lr){2-2}\cmidrule(lr){3-3}\cmidrule(lr){4-4}\cmidrule(lr){5-5}\cmidrule(lr){6-6}\cmidrule(lr){7-7}\cmidrule(lr){8-8}\cmidrule(lr){9-9}\cmidrule(lr){10-10}
& \multicolumn{1}{c}{(1)} & \multicolumn{1}{c}{(2)} & \multicolumn{1}{c}{(3)} & \multicolumn{1}{c}{(4)} & \multicolumn{1}{c}{(5)} & \multicolumn{1}{c}{(6)} & \multicolumn{1}{c}{(7)} & \multicolumn{1}{c}{(8)} & \multicolumn{1}{c}{(9)} \\ \toprule
\underline{\textit{Panel A: Pooled}}  \\ 
\textit{Intensity 1999 x post} &       -0.816*** &       -0.924*** &       -0.614*** &        0.149 &        0.792* &       -1.670*** &       -0.279 &        0.019 &        1.681*** \\ 
 & (       0.188) & (       0.183) & (       0.134) & (       0.281) & (       0.470) & (       0.210) & (       0.248) & (       0.048) & (       0.498) \\ 
\textit{RW p-value} & [ 0.000] & [ 0.000] & [ 0.000] & [ 1.000] & [ 0.370] & [ 0.000] & [ 0.784] & [ 1.000] & [ 0.004] \\ 
  & & & & & & & & & \\ 
Mean (pre-1997) &         2.48 &         2.58 &         0.91 &         4.86 &        15.84 &         2.56 &         3.54 &         0.27 &        14.39 \\ 
Obs &       15,642 &       15,642 &       15,642 &       15,642 &       15,642 &       15,642 &       15,642 &       15,642 &       15,642 \\ 
  & & & & & & & & & \\ 
\underline{\textit{Panel B: Females}}  \\ 
\textit{Intensity 1999 x post} &       -0.600*** &       -0.758*** &       -0.687*** &       -0.126 &        0.821 &       -1.792*** &       -0.447 &       -0.035 &        1.279** \\ 
 & (       0.215) & (       0.271) & (       0.170) & (       0.350) & (       0.577) & (       0.227) & (       0.308) & (       0.043) & (       0.579) \\ 
\textit{RW p-value} & [ 0.037] & [ 0.037] & [ 0.000] & [ 0.845] & [ 0.591] & [ 0.000] & [ 0.591] & [ 0.845] & [ 0.137] \\ 
  & & & & & & & & & \\ 
Mean (pre-1997) &         2.13 &         3.15 &         0.94 &         4.43 &        16.35 &         2.46 &         3.86 &         0.12 &        13.25 \\ 
Obs &       15,642 &       15,642 &       15,642 &       15,642 &       15,642 &       15,642 &       15,642 &       15,642 &       15,642 \\ 
  & & & & & & & & & \\ 
\underline{\textit{Panel C: Males}}  \\ 
\textit{Intensity 1999 x post} &       -1.074*** &       -1.086*** &       -0.559*** &        0.450 &        0.702 &       -1.532*** &       -0.140 &        0.075 &        2.163*** \\ 
 & (       0.262) & (       0.212) & (       0.159) & (       0.334) & (       0.565) & (       0.261) & (       0.263) & (       0.083) & (       0.641) \\ 
\textit{RW p-value} & [ 0.000] & [ 0.000] & [ 0.003] & [ 0.714] & [ 0.714] & [ 0.000] & [ 0.734] & [ 0.734] & [ 0.004] \\ 
  & & & & & & & & & \\ 
Mean (pre-1997) &         2.85 &         2.04 &         0.89 &         5.28 &        15.38 &         2.69 &         3.21 &         0.42 &        15.73 \\ 
Obs &       15,642 &       15,642 &       15,642 &       15,642 &       15,642 &       15,642 &       15,642 &       15,642 &       15,642 \\ 
  & & & & & & & & & \\ 
\bottomrule
\end{tabular}
    }
    \vspace{-3mm}
    \floatfoot{\deleted{Notes: Progresa intensity data are pooled from 1990 to 2000. All regression models include year fixed effects and municipality fixed effects. Standard errors are clustered at the municipality level. The regression model is weighted by the number of population aged +65. ***p$<0.01$, **p$<0.05$, *p$<0.1$.}\added{Notes: This table reports difference-in-differences estimates of the effect of Progresa enrollment intensity on cause-specific Elder Mortality Rate (EMR) for adults aged 65 and older per 1,000 inhabitants, using the Barham-Rowberry sample (municipalities incorporated into Progresa in 1998 or 1999). Each column corresponds to a specific cause of death. All regressions include municipality and year fixed effects. Standard errors are clustered at the municipality level. ***p$<0.01$, **p$<0.05$, *p$<0.1$.}}
  }
\end{table}
\end{landscape}


\begin{table}[htbp]
\centering
\caption{\deleted{Progresa Impacts on Household Expenditures: Heterogeneity by Elderly Presence (Eligible Households)}\added{Short-Run Effects of Progresa on Household Expenditure by Elderly Presence, Eligible Households}}
\label{tab:expenditures_elderly}
\floatbox[{\captop}]{table}[\FBwidth]
  {}
  {\scalebox{0.9}{
      \begin{tabular}{lcccc}
\toprule
 & (1) & (2) & (3) & (4) \\
 & ln Food & ln Medicine & \% Food & \% Medicine \\
\midrule
Treat $\times$ 1999 (No Elderly)   & 0.121$^{***}$ & 0.039 & $-$0.014 & 0.000 \\
                                    & (0.026) & (0.062) & (0.014) & (0.005) \\[4pt]
Diff: Elderly HH $\times$ 1999     & 0.007 & $-$0.012 & 0.021 & 0.004 \\
                                    & (0.035) & (0.115) & (0.017) & (0.011) \\[4pt]
Observations                        & 22,524 & 22,763 & 22,478 & 22,489 \\
Control Mean (1997)                 & 4.668 & 0.744 & 0.757 & 0.053 \\
\bottomrule
\end{tabular}



    }
    \vspace{-3mm}
    \floatfoot{\deleted{Notes: . ***p$<0.01$, **p$<0.05$, *p$<0.1$.}\added{Notes: This table reports difference-in-differences estimates of the effect of Progresa enrollment intensity on household expenditure outcomes, stratified by the presence of elderly members in the household. The sample is restricted to eligible households. All regressions include municipality and year fixed effects. Standard errors are clustered at the municipality level and shown in parentheses. ***p$<0.01$, **p$<0.05$, *p$<0.1$.}}
  }
\end{table}


\end{document}